\documentclass[11pt,a4paper]{article}
\usepackage[margin=25mm]{geometry}
\usepackage{amsmath,amssymb,amsthm,booktabs,hyperref}
\newtheorem{proposition}{Proposition}
\title{Identifying historical geometries from complete rosters}
\author{Human Compression project: bounded source correspondence}
\date{}
\begin{document}\maketitle
\section{A finite identification principle}
Fix a finite group $G$, a rank, and a class of flag-transitive geometries
with specified local hypotheses. A marked input is a family of actual
subgroups $(G_i)$. Two inputs are equivalent when one automorphism of $G$
maps the whole family to the other after a permutation of types. The
accepted shadow packets retain exactly this convention.

\begin{proposition}
Suppose a complete roster for this class is given, together with a
verified quotient into whole-family equivalence classes. Let $\Gamma$
be a historical geometry known to belong to the class. Suppose its
source gives the multiset $s$ of type sizes and the order $b$ of its
chamber stabilizer. If exactly one roster class has these two invariants,
then that class represents $\Gamma$ with its stated $G$-action, up to
an automorphism of $G$ and a permutation of types.
\end{proposition}
\begin{proof}
Choose a chamber in $\Gamma$ and let $G_i$ be its object stabilizers.
Flag transitivity identifies the objects of type $i$ equivariantly with
the right cosets of $G_i$; hence their number is $[G:G_i]$. The chamber
stabilizer is $\bigcap_iG_i$. A change of chamber conjugates the whole
family, an isomorphism of the acting group applies one automorphism to
it, and a change of type names permutes it. Therefore
\[
 \left(\operatorname{multiset}\{[G:G_i]\},
       \left|\bigcap_iG_i\right|\right)
\]
is invariant on each whole-family class. Completeness puts the historical
family into some roster class. It has the printed invariants, and the
verified uniqueness makes this the specified class. This proof uses
completeness as a premise; equality of aggregate numbers alone would not
supply it. Nor does equality of invariants merge any two roster entries:
the accepted exact transporters supply the family quotient.
\end{proof}

\begin{proposition}
Suppose the same complete roster has exactly one class at rank $r$.
Then every historical rank-$r$ geometry satisfying its hypotheses has
that class, regardless of its numerical name. Every isomorphism-invariant
property established for this class holds for each such historical geometry.
\end{proposition}
\begin{proof}
Apply completeness to the historical marked family. Its class must be
the only class in the roster. The type permutation and group automorphism
induce an incidence isomorphism, which transports every flag, residue,
shadow and intersection witness. Therefore each of $\mathrm{IP}_k$,
$\mathrm{WIP}_k$, full IP and full WIP has the same truth value.
This uses exhaustive coverage of the class, not merely agreement between
two reported counts.
\end{proof}

\section{The unique rank-five Mathieu geometry}
The retained complete phase-six $M_{12}$ census exhausts all 147 Borel
anchors and the larger class in which local two-transitivity is not
assumed. It has exactly one rank-five class, represented by witness 55.
Its Borel has order 3, its maximal parabolics have orders
$108,720,720,7920,7920$, and it satisfies local two-transitivity.
The exact prior source-to-census comparison identifies the sole record
of \texttt{m12.rk5} with witness 55. These are accepted computation
premises, including the complete subgroup catalogue and finite group
algorithms, not new computations or a new classification here.

Any $M_{12}$ rank-five geometry denoted $5.1$ in Tables 2 and 4 of
\cite{JL} belongs to this larger class. The preceding proposition therefore
identifies it without assuming that its suffix equals an archive position.
The accepted all-residue computation proves $\mathrm{IP}_3$ and
$\mathrm{IP}_4$ and disproves $\mathrm{WIP}_5$. Hence both named positive
claims hold, and the same geometry fails full WIP. All of these assertions
use the exact specified $M_{12}$ action and the retained Aut/type convention.

\section{The negative rank-six Higman--Sims assertion}
The unnumbered paragraph on printed page 289 of \cite{JL} names the
rank-six HS geometry $6.1$ and states that it fails $\mathrm{WIP}_6$.
This statement is separate from the 232 positive occurrences in its
Tables 2, 4, 6 and 8. The retained earlier HS classification has exactly
one rank-six class. The accepted quotient of the supplied roster represents
it by \texttt{hs.rk6}, active entry 1, class 80 (global input 124).
The singleton proposition identifies the historical geometry $6.1$ with
this family, under precisely that retained global completeness premise.
Its accepted all-residue record has $\mathrm{IP}_2$ but fails
$\mathrm{WIP}_k$ for each $k=3,4,5,6$; in particular it verifies the named
negative assertion. We neither infer completeness from the presence of
one archive entry nor treat the geometry's historical name as an input to
the shadow test. This adds one negative named assertion and does not
change the count of four settled positive occurrences or the 228
remaining positive correspondences.

\section{The Higman--Sims geometry numbered 3.2}
The local primary article of Leemans \cite{HS}, printed page 234, names
its $\mathrm{HoSi}.c^*$ construction as geometry $3.2$ in supplement [12].
Its displayed type sizes are $176,15400,1100$, and its Borel is $S_6$,
of order $720$. Reference [12] on page 236 identifies the same 2003 HS
supplement and author URL as reference [21] of Jacobs--Leemans
\cite{JL}. Thus the numeric name has an explicit documentary bridge;
it is not inferred from its position in the supplementary permutation
file.

The accepted HS computation processes all 124 active inputs and
quotients them into 80 actual whole-family classes. It retains the
historical classification as its completeness premise. Among all those
inputs, the displayed three sizes and Borel select precisely the two
rank-three active entries 2 and 3. Their subgroup orders are
$2880,252000,40320$, in that order, in the group of order $44352000$.
The accepted transporter identifies them as the same class, class 2
with representative active entry 2. No other class has this invariant.
By the proposition, under exactly that retained completeness premise,
this is the historical geometry $3.2$. Its accepted shadow computation
has the full intersection property. Consequently the two occurrences
of HS $3.2$ in Tables 2 and 8 of \cite{JL} have an identified positive
representative.

The second construction in \cite{HS}, pages 234--235, has type sizes
$176,15400,176$ and Borel order $720$. These select class 1 uniquely,
represented by active entry 1, whose accepted result is weak
intersection without full intersection. The article body does not print
its supplement number. We identify this \emph{named construction}, but
do not promote the tempting number $3.1$ to an established historical
label.

\section{Hall--Janko diagrams and the exact limit of their names}
The five-page primary article \cite{J2}, pages 79--80, prints three
rank-four diagrams as numbers 1, 2 and 3. Their invariants select the
three representatives of the accepted complete Hall--Janko census as
follows. The symbols $R_{28},R_{29},R_{30}$ refer to its exact witness
file, not to historical publication numbers.
\begin{center}\begin{tabular}{clcc}
\toprule
Article diagram & Multiset of type sizes & Borel order & Witness\\
\midrule
1 & $4032,4032,10080,10080$ & 1 & $R_{29}$\\
2 & $4032,10080,10080,50400$ & 1 & $R_{28}$\\
3 & $100,100,1575,1575$ & 8 & $R_{30}$\\
\bottomrule
\end{tabular}\end{center}
Every line follows by dividing $604800$ by the actual four subgroup
orders and comparing the complete three-class roster. Thus diagram 2
has full intersection, whereas diagrams 1 and 3 fail weak intersection.
No subgroup census or large shadow run is repeated here.

Jacobs--Leemans instead names a supplemental geometry $4.2$.
The already-local primary body does not explicitly identify its article
ordinal 2 with that supplemental identifier. The common rank-prefix
interpretation is plausible, but this packet keeps it as a documentary
obligation. In particular these three identifications do not discharge
any extra occurrence beyond the four settled claims among the 232 historical labels recorded in the
frozen coherent L096 assembly.

\section{What remains unresolved}
The current fixed-access register still archives the Janko $J_1$
classification chapter of Gottschalk--Leemans and the related
Magma-classification article as unavailable. The accessible thin-$J_1$
and Livingstone papers do not give the needed historical number $4.1$.
The existence of one positive and one negative rank-four class therefore
does not identify that name. No new retrieval was attempted.

The files called \texttt{catalogue\_L097.pdf},
\texttt{catalogue\_L098.pdf} and \texttt{catalogue\_L101.pdf} have exactly
the same hash and body as the Hall--Janko article \cite{J2}; they are not
primary sources for $J_3,M_{23},M_{22}$. We preserve this previously known
access limitation. The archive README explicitly identifies \texttt{hs.rk4}'s final
\texttt{geo[2]} as $4.1$ of a cited reference, whereas the accepted
calculation proves that exact input fails both $\mathrm{IP}_3$ and
$\mathrm{IP}_4$. Its literal reference number is [20], but the final
paper's HS supplement is [21]. Thus this is an internally inconsistent
edition-level routing assertion when combined with the positive table
claims. We record the conflict without deciding whether the example,
reference number, external label or table was misprinted. No other
label assignment follows from it.

This additive note identifies HS $3.2$ and the unique $M_{12}$ rank-five
class named $5.1$, proving four positive occurrences and leaving 228
positive-label occurrences without an established historical correspondence. It does not complete L096,
recover any additional original, or claim a new classification, theorem
priority or length factor. All accepted common proofs and complete
program/data dependencies remain paid; this finite invariant comparison
does not replace them.

\begin{thebibliography}{9}
\bibitem{HS} D.~Leemans, The residually weakly primitive geometries of HS,
Australas. J. Combin. 33 (2005), 231--236. Exact local primary:
\texttt{older\_L091.pdf}; identifiers read directly on printed pp.~234--235.
\bibitem{J2} D.~Leemans, The residually weakly primitive geometries of $J_2$,
Note Mat. 21 (2002), 77--81. Exact local primary:
\texttt{catalogue\_L104.pdf}; diagrams on printed pp.~79--80.
\bibitem{JL} P.~Jacobs and D.~Leemans, An algorithmic analysis of the
intersection property, LMS J. Comput. Math. 7 (2004), 284--299.
The source's Tables 2 and 8, and reference [21], are used here.
\end{thebibliography}
\end{document}
